\scanpage{81}
étale $\underline J$ de degré $3$ de $S$, permettant de définir un $\Sigma_J$

\[
\Sigma_J=\mathbb P(V_{\underline J}(\underline O_S))
\tag*{(25)}
\]

\diagram{\begin{tikzcd}[ampersand replacement=\&,column sep=small,row sep=small]
0\arrow[r,no head] \& V_J(O_S)\arrow[r,no head] \& \underline O_S^{\underline J}\arrow[r,"\inserted{\mathrm{tr}}"]\arrow[d,equal] \& \underline O_S\arrow[r] \& {}\\
\& \& p_*(\underline O_{\underline J})
\end{tikzcd}}

\deleted{\unsure{de|d}}\unsure{et|Et} un module inv. $\underline O_{\Sigma_J}(1)$, \inserted{où $p:\underline J\mathbin{\unsure{-|\to}}S$} au module

\[
T_{0\Sigma_J}=\underline O_{\Sigma_J}(1)^{\natural}
\tag*{(26)}
\]

de sorte qu'on a un iso

\[
T_{0\Sigma}^{\otimes2}\simeq \underline O_{\Sigma_J}(2)\wedge_{\{\pm1\}}\mathord{\mathrm{I}\mkern-7mu\mu}_4^*
\simeq T_{\Sigma_J}\deleted{\unsure{\otimes}}\wedge_{\{\pm1\}}(\mathord{\mathrm{I}\mkern-7mu\mu}_4^*\wedge\underline\omega_{\underline J})
\tag*{(27)}
\]

où $\underline\omega_J$ est le revêtement \inserted{étale} quadratique associé au revêtement \inserted{étale} cubique $\underline J$.

Si maintenant $E$ est une courbe elliptique sur la base $S$, on lui associe

\[
\underline J=({}_2E)^*\qquad\text{rev. étale cubique de }S
\tag*{(28)}
\]

et \deleted{on \unsure{a|en}} un iso

\[
E/(\pm\mathrm{id}_E)\simeq\Sigma_{\underline J}
\tag*{(29)}
\]

et un iso
% archive sheet number: 80
